\documentclass{article}
% these lines load any packages you might wish to use. 
% some common ones are loaded below
\usepackage{amsfonts, amsmath, amssymb}
\usepackage[numbers,square,sort&compress]{natbib}
\usepackage{graphicx}  % required for inserting images
\usepackage{hyperref}  % for inserting urls, and having clickable links to equation and figure numbers
\usepackage[margin=1.2 in]{geometry}  % these margins are easy to read
\geometry{letterpaper}                   % ... or a4paper or a5paper or ... 
\usepackage{lmodern}		% good font
\usepackage{caption} % for \captionsetup
\usepackage{appendix}
\usepackage{subcaption}
\usepackage{graphicx}
\usepackage{float}
\usepackage{hyperref}
\usepackage{gensymb}

\usepackage{tikz}  % for using tikz, to draw your own figures
\usetikzlibrary{angles,quotes}

% input document information here
% \title{Is the cook cooked?}
\title{Kitchen Safety: The danger of heat to the chef within a two-dimensional plane}
\author{Ryan Hsiao 61609715, Fanbo Feng 95719241, Davis Li 49877293}
\date{}  % if you uncomment this, it will not put the date.
%\date{Dec 99, 1999}   % if you uncomment this, it will list the date as whatever is inside the curly brackets
\begin{document}

\maketitle

\section{Introduction}
In professional and domestic kitchens, chefs are often exposed to high temperatures generated by cooking appliances, especially stoves. Prolonged exposure to such heat can cause thermal discomfort, reduced work efficiency, and even heat-related illnesses. The thermal environment within the kitchen depends on various factors, including the intensity of the heat source and the effectiveness of cooling systems such as air conditioning. Since heat transfer involves diffusion, sources, and dissipation processes, mathematical models provide a useful framework for studying such thermal environments.

\vspace{3mm}

\noindent To analyze this problem, we model the kitchen as a two-dimensional square domain containing a fixed stove that continuously releases radiant heat laterally into the air. We simulate how temperature evolves over time using a partial differential equation (PDE) that combines heat diffusion, a localized heat source, and heat dissipation due to air conditioning. Using numerical simulations based on the Backward Euler scheme, we generate solutions for different air-conditioning strengths. The resulting temperature distributions are then analyzed using heatmaps and time-series plots in order to interpret the characteristics of the thermal environment. By comparing the steady-state temperatures and the size of dangerous hot zones, we aim to identify how air-conditioning strength influences the safety and comfort of kitchen workspace.
%%%%%%%%%%%%%%%%%%%%% Model Section Start %%%%%%%%%%%%%%%%%%%%%
\section{Methodology}\label{sec:model}

\subsection{Definition and Notations}

%%%%%%%%%%%%%%%%%%%%%%%%%%%%%%%%%%%%%%%%%% Table start
\begin{table}[!h]
\begin{center}
\begin{tabular}{||c | p{7.5cm} | c | c||} 
 \hline
 Symbol & Description & Type & Dimensions \\ [0.3ex] 
 \hline\hline
 $u$ & Temperature at location (x, y) and time t. & Dependent Variable & $\theta$ \\ 
 \hline
 $x$ & Horizontal position in the space. & Independent Variable & L \\ 
 \hline
 $y$ & Vertical position in the space.  & Independent Variable & L \\ 
 \hline
 $z$ & Height of the space. & parameter & L \\ 
 \hline
 $t$ & Time. & Independent variable & T \\ 
 \hline
 $T_a$ & Room temperature. & Parameter & $\theta$ \\ 
 \hline
 $\alpha$ & Thermal diffusivity of air. & Parameter & L$^2$T$^{-1}$ \\ 
 \hline
 $\beta$ & Heat dissipation rate due to air conditioning (AC). & Parameter & T$^{-1}$ \\ 
 \hline
 $\gamma$ & Radiant power emitted laterally by the stove into the air. & Parameter & $\theta$ L$^3$T$^{-1}$ \\ 
 \hline
 $\Omega$ & The domain of the space. & Set & - \\ 
 \hline
\end{tabular} 
\end{center}
\captionsetup{font=small}
\caption{The variables and parameters are measured in \textit{hours} ($h$, dimension T), \textit{meters} ($m$, dimension L), and \textit{Kelvin} ($K$, dimension $\theta$).} % All dimensions L has the unit of meter, and dimensions of T has the unit of second.
\end{table}
%%%%%%%%%%%%%%%%%%%%%%%%%%%%%%%%%%%%%%%%%% Table end

\noindent We consider a kitchen with a stove in a fixed location that constantly releases radiant heat into the air laterally. This is modeled in a two-dimensional domain as we ignore the effect of the height of the kitchen. 

\vspace{3mm}

\noindent The purpose of this model is to analyze how heat from a stove affects the temperature experienced by a cook and the surrounding thermal environment. The domain of the room $\Omega$ represents the available space that heat can spread, which we assume to be square-shaped. We determine a set parameter for the area $\Omega$ by establishing a coordinate system on the kitchen. We define the lower-left corner of the space as the origin, with the bottom wall representing the $x$-axis and the left wall representing the $y$-axis. We define the region as follows:
\begin{align*}
    \Omega = \left\{ (x, y)\in\mathbb{R}^2 \mid 0\le x \le 3 \ , 0 \le y \le 3 \right\} \ , 
\end{align*}

\noindent where the stove is located at the point $(2.8, 1.5)$.

\vspace{3mm}

\noindent Next, we assume conditions to simplify the model from a real-world scenario in order to achieve the target model mentioned above. First, the model takes place in a two-dimensional indoor setting. The height of the kitchen is ignored so that we average out the temperature in the direction $z$. Furthermore, the only method to remove heat from the space is by central air conditioner---the walls do not absorb or emit heat. The central air conditioner will reduce the same amount of heat across the kitchen. Next, the stove is the only heat source and can provide heat continuously to simulate a continuous cooking environment. In addition, the cookware does not affect the heat transfer from the stove to the air. Although the diffusivity of the air varies with temperature, we assume it to be constant in our model. These assumptions allow us to capture heat conduction over $\Omega$. 

\subsection{Model Formulation}
In our model, heat is introduced into the system with a source and is removed with air conditioning. The heat introduced by the stove is modeled by a function $S(x, y)$, which represents the spatial distribution of the heat source. We assume that the total heat input satisfies
\begin{align*}
z\displaystyle\iint S(x, y) dxdy = \gamma. 
\end{align*}
Following [4], we model $S(x, y)$ as a standardized two-dimensional Gaussian distribution:
\begin{align*}
    S(x, y) = \frac{\gamma/z}{2\pi\sigma^2}\exp\left(-\frac{(x-2.8)^2 + (y-1.5)^2}{2\sigma^2}\right) \ ,
\end{align*}

\noindent where $\sigma$ is the standard deviation of the Gaussian and $\gamma > 0$ denotes the radiant power emission rate of the stove. 

\vspace{3mm}

\noindent Next, heat loss due to air conditioning, or the sink term, can be defined using Newton’s law of cooling,
\begin{align*}
\dfrac{dT}{dt} = -\beta(T - T_a) \, 
\end{align*}
where T represents the temperature and is replaced by $u(x,y,t)$ in our model: 
\begin{align*}
    L(u) = -\beta(u - T_a) \ ,
\end{align*}

\noindent where $\beta > 0$ denotes the heat dissipation rate due to AC. 

\vspace{3mm}

\noindent Finally, we have a two dimensional diffusion equation 
\begin{align*}
    u_t = \alpha(u_{xx} + u_{yy}) \ ,
\end{align*}

\noindent where $\alpha>0$ describes the thermal diffusivity of air, and $u_{xx}$ and $u_{yy}$ denote the lateral change in temperature in the $x$ and $y$ directions, respectively.

\vspace{3mm}

\noindent Combining the source and sink together with the diffusion equation, we get the following 
\begin{align*}
    u_t = \alpha(u_{xx} + u_{yy}) + S(x, y) + L(u) \ , \tag{2.1a}
\end{align*}

\noindent with initial condition 
\begin{align*}
    u(x, y, 0) = T_a \ .\tag{2.1b}
\end{align*}

\noindent The insulation of the walls gives the boundary conditions as 
\begin{align*}
    u_x(0, y, t) &= u_x(3, y, t) = 0 \ , \tag{2.1c} \\ 
    u_y(x, 0, t) &= u_y(x, 3, t) = 0 \ . \tag{2.1d}
\end{align*}

\noindent Together, equations (2.1a) to (2.1d) form a PDE system governing the temperature distribution in the kitchen.

\subsection{Numerical Simulation}
\noindent In this section, we approximate the solution to the PDE system (2.1) by backward Euler (BE) scheme in time and central differences in space. Since we perform simulation using numerical methods, the step sizes $\Delta t$, $\Delta x$, and $\Delta y$ need to be minimal in order to obtain numerical stability. Hence, we let $\Delta x = \Delta y = h$ and $\Delta t = 0.1 \cdot h$, where $h = \dfrac{3}{J}$ with $J = 150$ denoting the number of spatial grid points along each axis. In our context, stability indicates that the solution does not grow exponentially at any point.

\subsubsection*{Discretization of the PDE}
\noindent Applying the BE scheme in time and central differences in space to (2.1a) to get the following with two-dimensional grid $(i, j)$: 
\begin{align*}
    \frac{u^{n+1}_{i, j} - u^{n}_{i,j}}{k} = \alpha \left( \frac{u^{n+1}_{i-1, j} - 2u^{n+1}_{i, j} + u^{n+1}_{i+1, j}}{h^2} + \frac{u^{n+1}_{i, j-1} - 2u^{n+1}_{i, j} + u^{n+1}_{i, j+1}}{h^2} \right) + S(x, y) - \beta u^{n+1}_{i, j} + \beta T_a \ , \tag{2.2}
\end{align*}

\noindent where $\Delta t = k$ and $\Delta x = \Delta y = h$ as the domain is a square, and $- \beta u^{n+1}_{i, j} + \beta T_a$ is expanded from $L(u)$.

\vspace{3mm}

\noindent Multiplying both sides of (2.2) by $k$ and rearranging, we let $r = \dfrac{\alpha k}{h^2}$ to obtain 
\begin{align*}
    u^n_{i, j} &= u^{n+1}_{i, j} - r \left( u^{n+1}_{i-1, j} - 2u^{n+1}_{i, j} + u^{n+1}_{i+1, j} + u^{n+1}_{i, j-1} - 2u^{n+1}_{i, j} + u^{n+1}_{i, j+1} \right) - kS(x, y) + k\beta u^{n+1}_{i, j} - k\beta T_a \\ 
    &= u^{n+1}_{i, j} - ru^{n+1}_{i-1, j} + 2ru^{n+1}_{i, j} -r u^{n+1}_{i+1, j} -r u^{n+1}_{i, j-1} +2ru^{n+1}_{i, j} -r u^{n+1}_{i, j+1} - kS(x, y) + k\beta u^{n+1}_{i, j} - k\beta T_a \\
    &= (1 + 4r + k\beta) u^{n+1}_{i, j} - r u^{n+1}_{i-1, j} - r u^{n+1}_{i+1, j} - r u^{n+1}_{i, j-1} - r u^{n+1}_{i, j+1} - kS(x, y) - k\beta T_a \ . \tag{2.3}
\end{align*}

\noindent Now, we can express (2.3) in matrix form as
\begin{align*}
    & \textbf{u}^n = A \textbf{u}^{n+1} - kS(x, y) - k\beta T_a  \\
    \implies & A \textbf{u}^{n+1} =\textbf{u}^n + kS(x, y) + k\beta T_a \\
    \implies & A \textbf{u}^{n+1} = \textbf{b} \ , \tag{2.4}
\end{align*}

\noindent where $A = I - r L + k\beta I$ is a sparse matrix with $L$ being the unscaled standard five-points discrete Laplacian obtained by the Kronecker product, $I$ is the identity matrix, and $\textbf{b} = \textbf{u}^n + kS(x, y) + k\beta T_a$. In addition, we add Neumann boundary conditions (2.1c) and (2.1d) to the system by changing the boundary in matrix $L$.

\subsubsection*{Simulation Process}
\noindent From the above section, we transformed the discretized PDE (2.2) into a linear algebra problem (2.4). At each time step, we solve the linear system (2.4) using the \texttt{spsolve} function from \texttt{scipy.sparse.linalg} in Python, which employs a direct LU factorization. The iteration terminates once the maximum number of time steps is reached. Finally, we get a result of heat distribution over the domain $\Omega$. The code is available in the \href{https://github.com/tusm06/MATH461_Unit03}{Git Repository}.
%%%%%%%%%%%%%%%%%%%%% Model Section End %%%%%%%%%%%%%%%%%%%%%


%%%%%%%%%%%%%%%%%%%%% Analysis Section Start %%%%%%%%%%%%%%%%%%%%%
\section{Analysis of the model}

\subsection{Stability of the PDE}

In this section, we solve the PDE and compute its general solution, which gives the temperature function $u(x,y,t)$. To achieve this goal, we need use separation of variables and calculate the solution in two separate parts. Based on the result we get, we can analyze the steady-state of our equation.

\vspace{3mm}

\noindent Firstly, we substitute $L(u)$ into our equation (2.1a), and rewrite it in the form of following:
\begin{align*}
    u_t - \alpha(u_{xx} + u_{yy}) + \beta u = A (x, y) \ , \tag{3.1}
\end{align*}
where $A(x,y) = S(x, y) + \beta T_a$, represents the combined source term. Subsequently, we find a component term of $u (x,y,t)$ that is only in the respect of $x$ and $y$. We define it as $w(x,y) = X(x)Y(y)$. Based on the homogeneous version of equation 3.1 and letting the right hand side of the equation to be zero, we get the equation
$$
w_{xx} + w_{yy} + \lambda w = 0 \ .
$$
Hence, we can compute the eigenfunctions $w(x,y)$ and eigenvalues $\lambda$ by applying the boundary condition 2.1c and 2.1d as following
\begin{align*}
    w_{mn}(x,y) &= \cos\left(\frac{m\pi x}{3}\right)\cos\left(\frac{n\pi y}{3}\right), \quad m,n = 0,1,2,... \\
    \lambda_{mn} &= \left(\frac{m\pi}{3}\right)^2 + \left(\frac{n\pi}{3}\right)^2, \quad m,n = 0,1,2,... \ .
\end{align*}
Then, we introduce a new term $v(t)$ write $u(x,y,t)$ as
\begin{align*}
    u(x,y,t) = \sum_{m}^\infty \sum_{n}^\infty w_{mn}(x,y)v_{mn}(t) \ . \tag{3.2}
\end{align*}

\vspace{3mm}

\noindent In order to compute another term in our solution, we expand the amplitude term $A(x,y)$ as a cosine function in terms of $w_{mn}(x,y)$. We define the coefficient for this term to be $a_{mn}$, so we have $A(x,y)= \sum_{m,n} a_{mn}w_{mn}(x,y)$. Therefore, we substitute the expanded amplitude term $A(x,y)$ in equation 3.1, and break the PDE into an ODE in the respective of $t$ and solve it to obtain
\begin{align*}
    v_{mn}(t) = b_{mn}e^{-(\alpha \lambda_{mn}+\beta)t}+\frac{a_{mn}}{\alpha \lambda_{mn}+\beta} \ , \tag{3.3}
\end{align*}
where $b_{mn}$ is our newly defined coefficient. Then, apply our initial condition (2.1b), we get $u(x,y,0) = v_{mn}(0)w_{mn}(x,y)$. Substituting it back into equation (3.3), we get the expression of $b_{mn}$ as following:
$$
b_{mn} = v_0 - \frac{a_{mn}}{\alpha \lambda_{mn}+\beta} \ .
$$
Hence, we get the solution of $v_{mn}$ as below:
\begin{align*}
    v_{mn}(t) = \left(v_0 - \frac{a_{mn}}{\alpha \lambda_{mn}+\beta}\right)e^{-(\alpha \lambda_{mn}+\beta)t}+\frac{a_{mn}}{\alpha \lambda_{mn}+\beta} \ .
\end{align*}

\noindent Now, we can get the full general solution by substituting $w_{mn}(x,y)$ and $v_{mn}(t)$ back into equation (3.2):
\begin{align*}
    u(x,y,t) =  u_{steady}(x,y)+ \sum_{m}^\infty \sum_{n}^\infty \left(v_0 - \frac{a_{mn}}{\alpha \lambda_{mn}+\beta}\right)e^{-(\alpha \lambda_{mn}+\beta)t}\cos\left(\frac{m\pi x}{3}\right)\cos\left(\frac{n\pi y}{3}\right) \ ,  \tag{3.4}
\end{align*}
where $u_{steady}(x,y)$ denotes the steady-state temperature distribution, given by
$$
u_{steady}(x,y) = \sum_{m}^\infty \sum_{n}^\infty  \frac{a_{mn}}{\alpha \lambda_{mn}+\beta}\cos\left(\frac{m\pi x}{3}\right)\cos\left(\frac{n\pi y}{3}\right) \ .
$$

\noindent Furthermore, equation (3.4) contains a single time-dependent term, $e^{-(\alpha \lambda_{mn}+\beta)t}$, on the right-hand side. As $t \to \infty$, this term converges to 0, so all terms except $u_{steady}(x,y)$ will vanish. Therefore, the solution approaches the steady-state distribution $u_{steady}(x,y)$ as $t\to \infty$.

\subsection{Analysis of the Numerical Solution}
\noindent From the previous section, we showed that the PDE system (2.1) becomes stable as $t\rightarrow\infty$. In this section, we analyze the numerical stability of the simulation. In our numerical simulation, we discretize the PDE system (2.1) into (2.3) using the Method of Lines (MOL) with backward Euler for time and central differences in space. We choose the backward Euler method because it is unconditionally stable. However, the presence of the source term $S(x, y)$ and sink term $L(u)$ may influence the numerical behavior of the scheme. Thus, we perform a von Neumann stability analysis to verify the stability of the numerical scheme and examine the local truncation error to establish consistency.

\subsubsection*{Von Neumann Stability Analysis}
\noindent We use Ansatz method to perform the analysis. Specifically, we assume 
\begin{align*}
u^n_{i, j} = g(\xi_i, \xi_j)\exp(I\xi_ihi + I\xi_jhj) \ , 
\end{align*}
where $x_i = ih$ and $y_j = jh$ denote the coordinates of the grid point $(i, j)$, $\xi_i$ and $\xi_j$ are wave numbers, and $g(\xi_i, \xi_j)$ is the corresponding amplification factor. Here, $I$ denotes the imaginary unit since $i$ is used for indexing. Substituting this into the second equation in (2.3) yields
\begin{align*}
g(\xi_i, \xi_j)^n &\exp(I\xi_i hi + I\xi_j hj) = g(\xi_i, \xi_j)^{n+1} \exp(I\xi_i hi + I\xi_j hj) \\
&- rg(\xi_i, \xi_j)^{n+1}(\exp(I\xi_i h(i-1) + I\xi_j hj) - 2\exp(I\xi_i hi + I\xi_j hj) + \exp(I\xi_i h(i+1) + I\xi_j hj) \\
&+ \exp(I\xi_i hi + I\xi_j h(j-1)) - 2\exp(I\xi_i hi + I\xi_j hj) + \exp(I\xi_i hi + I\xi_j h(j+1))) \\
& + k\beta g(\xi_i, \xi_j)^{n+1} \exp(I\xi_i hi + I\xi_j hj) \ .
\end{align*}

\noindent Dividing both sides by $g(\xi_i, \xi_j)^n\exp(I\xi_i hi + I\xi_j hj)$, we obtain 
\begin{align*}
   1 &= g(\xi_i, \xi_j) - rg(\xi_i, \xi_j)(\exp(-I\xi_i h) - 2 + \exp(I\xi_i h) + \exp(-I\xi_j h) - 2 + \exp(I\xi_j h)) + k\beta g(\xi_i, \xi_j) \\
    &= g(\xi_i, \xi_j) (1 - r(\exp(-I\xi_i h) - 2 + \exp(I\xi_i h) + \exp(-I\xi_j h) - 2 + \exp(I\xi_j h)) + k\beta) \ .
\end{align*}

\noindent Based on Euler's equation $\exp(I\theta) = \cos(\theta) + I\sin(\theta)$ and $\exp(-I\theta) = \cos(\theta) - I\sin(\theta)$, we get 
\begin{align*}
&1 = g(\xi_i, \xi_j) (1 - r(2\cos(\xi_i h) - 2 + 2\cos(\xi_j h) - 2) + k\beta) \\
\implies & g(\xi_i, \xi_j) = \frac{1}{1 - r(2\cos(\xi_i h) - 2 + 2\cos(\xi_j h) - 2) + k\beta} \\
\implies & g(\xi_i, \xi_j) = \frac{1}{D} \ , \tag{3.5}
\end{align*}

\noindent where $D = 1 - r(2\cos(\xi_i h) - 2 + 2\cos(\xi_j h) - 2) + k\beta$. 

\vspace{3mm}

\noindent For the numerical scheme to be stable, we require $|g(\xi_i, \xi_j)| \le 1$: we need $D \ge 1$. Considering only the $\cos(\xi_i h)$ term, we obtain the following bounds:
\begin{align*}
    -1 \le \cos&(\xi_i h) \le 1 , \\ 
    -2 \le 2\cos&(\xi_i h) \le 2 , \\
    -4 \le 2\cos&(\xi_i h) - 2 \le 0 .
\end{align*}

\noindent We conclude that the term $2\cos(\xi_ih)-2$ ranges between $[-4, 0]$, and the same holds for $2\cos(\xi_j h)-2$. Next, we know that $r = \dfrac{\alpha k}{h^2} > 0$ since $\alpha, k, h > 0$ as stated in the model formulation section. Finally, we also get $k\beta>0$ since $k, \beta>0$. Combining these results, we consider two extreme cases: 
\begin{enumerate}
    \item If both $2\cos(\xi_i h) - 2$ are at their lowest point $-4$, then 
\begin{align*}
    D = 1 - r(-4 - 4) + k\beta = 1 + 8r + k\beta \ge 1 \ ,
\end{align*}

\item If both $2\cos(\xi_i h) - 2$ are at their highest point $0$, then
\begin{align*}
    D = 1 - r(0) + k\beta = 1 + k\beta \ge 1 \ .
\end{align*}
\end{enumerate}

\noindent Hence, we conclude that $D$ is always greater than or equal to 1 as required. Therefore, $|g(\xi_i, \xi_j)| = \left| \dfrac{1}{D} \right| \le 1$ as required. We have shown that $|g(\xi_i, \xi_j)|\le 1$ holds for all wave numbers $\xi_i$ and $\xi_j$, independent of $k$ and $h$, the BE scheme is unconditionally stable.

\subsubsection*{Numerical Consistency}
\noindent To ensure numerical consistency, we examine the local truncation errors in both time and space. The backward Euler method is first-order in time as shown in the left hand side of equation (2.2), giving a temporal error $O(k)$.Central differences are second-order in space, yielding a spatial error of $O(h^2)$ in both $x$ and $y$, where $\Delta x = \Delta y = h$. Moreover, the term $L(u)$is linear and does not involve spatial derivatives, so the term $\beta u$ does not contribute to the spatial truncation error.

\vspace{3mm}

\noindent Hence, we have the local truncation error (LTE) $\tau = O(k) + O(h^2)$. We know the numerical scheme is consistent since $\tau \rightarrow0$ as $k\rightarrow0$ and $h\rightarrow0$. Hence, the numerical scheme is consistent.

\subsubsection*{Numerical Convergence}
By the Fundamental theorem of convergence, consistency and stability together imply convergence. Since we have shown both properties are satisfied above, the numerical solution converges to the true solution as $k\rightarrow0$ and $h\rightarrow0$.
%%%%%%%%%%%%%%%%%%%%% Analysis Section End %%%%%%%%%%%%%%%%%%%%%


\section{Results}\label{sec:Results}

\subsection{Assigning values for parameters}
\noindent We take the thermal diffusivity of air as $\alpha = 0.06984\ m^2/h$ based on material property databases, while we set the air density $\rho = 1.29\ kg/m^3$ and specific heat capacity $c_p = 1004\ J/kg\ \cdotp K$, which are standard values for air at room conditions used in heat transfer modeling [1]. We choose the optimal room temperature $T_a = 23 \degree$C because it is the median within the recommended indoor comfort range of $20-26 \degree$C given by CSA workplace ergonomics standards [2]. These values ensure that our model reflects realistic physical properties of air and typical indoor conditions.

\vspace{3mm}

\noindent The heat source parameter $\gamma$ is estimated from a experimental studies of a stove heat transfer. Studies on stove performance indicate a model of the heat transfer in a stove system, and the heat loss of the chamber wall is $171 \ W$ under a $20 \degree$C ambient indoor temperature and a stove with $52 \ kW$ power [3]. In our model, we only need to consider the heat transferred from the stove area to the entire kitchen room, which is only take a small part of the whole stove thermal power. Therefore, we consider the heat that emitted from the stove chamber as the heat source in our model, hence we set $Q = 0.171 \ kW$.
% \noindent The heat source parameter $\gamma$ is estimated from experimental studies of cooking oil fires and stove heat transfer. Experimental measurements report that an electric stove used to heat cooking oil have a thermal power of approximately $3\ kW$ under ambient indoor conditions ($20 \degree$C and 50\% relative humidity) [3]. Studies on stove performance also indicate that radiation accounts for approximately 47.16\% of total heat transfer [4]. Using this radiative fraction, the effective radiant heat input can be estimated as 
% \begin{align*}
%     Q_{rad} = fraction\ \cdotp\ Q_{release} = 0.4716(3 \ kW) \approx 1.4148\ kW \ .
% \end{align*}
To incorporate this into our temperature-based PDE model, we convert energy input into temperature, establishing the relation 
\begin{align*}
    Q = \rho c_p V \frac{\mathrm{d}T}{\mathrm{d}t} \ \Rightarrow \gamma = \frac{Q}{\rho c_p} \ , \tag{4.1}
\end{align*}
where $Q$ denotes the heat input rate (power) in watts, $V$ is the volume of air being heated, and $T$ stands for air temperature. This relation follows from an energy balance for the air, where the heat input is equated to the rate of change of thermal energy in the volume. Substituting the parameter values yields 
\begin{align*}
    \gamma = \frac{Q}{\rho c_p} = \frac{171\ W}{1.29\ kg/m^3(1004\ J/kg\ \cdotp K)} \approx 0.132\ K \cdotp m^3/s \approx  475.3 \ K \cdotp m^3/h
\end{align*}
This expresses the heat source as a spatially distributed temperature input consistent with our governing heat equation. In our model, heat loss is assumed to occur only through air conditioning, so the heat dissipation rate $\beta$ is treated as a tunable parameter representing the strength of the cooling system. Therefore, in our simulations, all parameters are fixed except for $\beta$, which is varied to study the effect of air conditioning strength on the temperature distribution.

\subsection{Visualizing the results}
After building the the model with different values of $\beta$, we visualize our results in the forms of 2D diffusion heatmap and temperature versus time plot. The heat maps show how the heat released by the stove spreads across the kitchen, while the time plot shows how quickly the temperature experienced by the cook approaches its steady-state value. Therefore, these plots provide a clear illustration of how the temperature evolves over the 12-hour time range for different air-conditioning strengths.


%%%%%%%%%%%%%%%%%%%%%%%%%%%%%%%% Begin plots %%%%%%%%%%%%%%%%%%%%%%%%%%%%%%%%%%%%%
\begin{figure}[H]
     \centering
     % --- First Row ---
     \begin{subfigure}[b]{0.3\textwidth}
         \centering
         \includegraphics[width=\textwidth]{images/output11.png}
          \captionsetup{font=small}
         \caption{The temperature of the room at steady-state with $\beta = 1$}
     \end{subfigure}
     \hfill
     \begin{subfigure}[b]{0.3\textwidth}
         \centering
         \includegraphics[width=\textwidth]{images/output21.png}
         \captionsetup{font=small}
         \caption{The temperature of the room at steady-state with $\beta = 3.51$}
     \end{subfigure}
     \hfill
     \begin{subfigure}[b]{0.3\textwidth}
         \centering
         \includegraphics[width=\textwidth]{images/output31.png}
         \captionsetup{font=small}
         \caption{The temperature of the room at steady-state with $\beta = 6$}
     \end{subfigure}
     \captionsetup{font=small}
     \caption{Heat maps of kitchen. The small red circle represents the position of stove, the larger red circle represents the area which will touch the chef, the green circle is representing the chef, and black point is the chef position we focused. Warmer colors indicate higher temperature, and vice-versa.}
     \label{fig:four_plots}
\end{figure}
%%%%%%%%%%%%%%%%%%%%%%%%%%%%%%%% End plots %%%%%%%%%%%%%%%%%%%%%%%%%%%%%%%%%%%%%

\subsubsection{The kitchen with small power AC}

We first consider the case of weak air conditioning, corresponding to a small heat dissipation rate. To make the result present the significant observation, we set the $\beta = 1$, with the unit $hour^{-1}$. 

\vspace{3mm}

\noindent In this case, the heat map in Figure 1 (a) shows that a large high-temperature region forms around the stove and remains significant throughout the simulation. As time elapses, the heat diffuses outward from the stove location, and the dangerous hot region occupies a relatively large portion of the kitchen. 

\vspace{3mm}

% %%%%%%%%%%%%%%%%%%%%%%%%%%%%%%%% Begin plots %%%%%%%%%%%%%%%%%%%%%%%%%%%%%%%%%%%%%
% \begin{figure}[h]
%     \centering
%     \includegraphics[width=0.8\textwidth]{images/output11.png}
%     \captionsetup{font=small}
%     \caption{Heat map with 4-hours interval during a 12 hours cooking process with choice of $\beta = 1$. The small red circle represents the position of stove, the larger red circle represents the area which will touch the chef, the green small circle is the position of chef, and black point represents the point that chef feel the temperature.}
% \end{figure}
% %%%%%%%%%%%%%%%%%%%%%%%%%%%%%%%% End plots %%%%%%%%%%%%%%%%%%%%%%%%%%%%%%%%%%%%%

\noindent From Figure 2 we can observe that the temperature at the cook position continuously rises from the ambient temperature and approaches a high steady-state value around the 5 hours. In addition, the green horizontal line in the graph is represented the safe temperature a human can endure, which we predetermine it as 313.15 K(40 $\degree$C). Clearly, the final temperature is far above the safety threshold, indicating that weak cooling is insufficiefnt to maintain a safe working environment near the stove. With low AC power, heat accumulates strongly and exposes the chef to severe thermal risk.


%%%%%%%%%%%%%%%%%%%%%%%%%%%%%%%% Begin plots %%%%%%%%%%%%%%%%%%%%%%%%%%%%%%%%%%%%%
\begin{figure}[h]
    \centering
    \includegraphics[width=0.6\textwidth]{images/output12.png}
    \captionsetup{font=small}
    \caption{Temperature at the cook position over time for $\beta = 1$. The blue, green, and red lines show the ambient temperature, safety threshold, and the steady-state temperature reached during the 12-hour simulation respectively. The blue dots shows the simulation results and form a line over time. The yellow stars indicate the snapshots position for the corresponding heatmap.}
\end{figure}
%%%%%%%%%%%%%%%%%%%%%%%%%%%%%%%% End plots %%%%%%%%%%%%%%%%%%%%%%%%%%%%%%%%%%%%%


\subsubsection{The kitchen with medium power AC}

Next, we consider the case of medium air conditioning power, corresponding to a moderate heat dissipation rate. In this case, we set $\beta = 3.51$ with the unit $hour^{-1}$. We choose this value because we want to find the particular value for $\beta$ by tuning the values of $\beta$ such that it can make the temperature concentrate at the safety temperature when it approach to the steady-state. As a consequence, we choose the approximately value 3.51 to achieve our goal. 

\vspace{3mm}

\noindent In figure 1 (b), the heat map shows that the high-temperature region around the stove is still visible, but the heated region becomes smaller and more concentrated near the stove location. This indicates that the medium power AC can remove heat more effectively and prevent the heat from spreading to a larger portion of the kitchen.

\vspace{3mm}

% %%%%%%%%%%%%%%%%%%%%%%%%%%%%%%%% Begin plots %%%%%%%%%%%%%%%%%%%%%%%%%%%%%%%%%%%%%
% \begin{figure}[h]
%     \centering
%     \includegraphics[width=0.8\textwidth]{images/output21.png}
%     \captionsetup{font=small}
%     \caption{Heat map with 4-hours interval during a 12 hours cooking process with choice of $\beta = 3.51$. The small red circle represents the position of stove, the larger red circle represents the area which will touch the chef, the green small circle is the position of chef, and black point represents the point that chef feel the temperature.}
% \end{figure}
% %%%%%%%%%%%%%%%%%%%%%%%%%%%%%%%% End plots %%%%%%%%%%%%%%%%%%%%%%%%%%%%%%%%%%%%%

\noindent In more detail, from Figure 3 we can observe that the temperature at the cook position still rises rapidly from the ambient temperature, but it approaches a much lower steady-state value compared with the case in the last section. In addition, the final temperature is very close to the safety threshold. This means that the medium power AC significantly reduces the thermal risk to the chef, although the cook is still exposed to a temperature near the safety limit. Therefore, the medium power AC provides a more balanced situation, which can be considered as a threshold situation across all choice of $\beta$.

%%%%%%%%%%%%%%%%%%%%%%%%%%%%%%%% Begin plots %%%%%%%%%%%%%%%%%%%%%%%%%%%%%%%%%%%%%
\begin{figure}[h]
    \centering
    \includegraphics[width=0.6\textwidth]{images/output22.png}
    \captionsetup{font=small}
    \caption{Temperature at the cook position over time for $\beta = 3.51$. The blue, green, and red lines show the ambient temperature, safety threshold, and the steady-state temperature reached during the 12-hour simulation respectively. The blue dots shows the simulation results and form a line over time. The yellow stars indicate the snapshots position for the corresponding heatmap.}
\end{figure}
%%%%%%%%%%%%%%%%%%%%%%%%%%%%%%%% End plots %%%%%%%%%%%%%%%%%%%%%%%%%%%%%%%%%%%%%

\subsubsection{The kitchen with large power AC}

Finally, we consider the case of strong air conditioning, corresponding to a large heat dissipation rate. We set $\beta = 6$ with unit $hour^{-1}$. 

\vspace{3mm}

\noindent In the Figure 1 (c), the temperature increase is highly confined to a small region around the stove, while most of the kitchen remains close to room temperature during the simulation.

\noindent When we look into the cook position plots represented as Figure 4, the temperature at the cook position increases only slightly from the ambient temperature and quickly stabilizes at a value below the safety threshold. This indicates that strong cooling can effectively control the heat experienced by the chef. Compared with the previous two cases, this case has the smallest heated region and the lowest steady-state temperature at the cook position. Hence, among the three cases, the large-power AC gives the safest kitchen environment.

%%%%%%%%%%%%%%%%%%%%%%%%%%%%%%%% Begin plots %%%%%%%%%%%%%%%%%%%%%%%%%%%%%%%%%%%%%
\begin{figure}[h]
    \centering
    \includegraphics[width=0.6\textwidth]{images/output32.png}
    \captionsetup{font=small}
    \caption{Temperature at the cook position over time for $\beta = 6$. The blue, green, and red lines show the ambient temperature, safety threshold, and the steady-state temperature reached during the 12-hour simulation respectively. The blue dots shows the simulation results and form a line over time. The yellow stars indicate the snapshots position for the corresponding heatmap.}
\end{figure}
%%%%%%%%%%%%%%%%%%%%%%%%%%%%%%%% End plots %%%%%%%%%%%%%%%%%%%%%%%%%%%%%%%%%%%%%

\section{Discussion}\label{sec:Discussion}

The results computed by 3 different situations show that the heat dissipation rate $\beta$ has a strong effect on the kitchen temperature. By comparing the three temperature plots, we see that as we choose the larger value for $\beta$, the lower final temperature will be computed, hence faster approach the stable state. The same pattern can also be found in heat map, the large $\beta$ will result in the smaller dangerous heated zone. Therefore, we can say that increasing $\beta$ reduces both the temperature level and the size of the dangerous heat zone, and hence lead to a safer kitchen environment.

\vspace{3mm}

\noindent In addition, all three cases approach a steady state over time, which agrees with the analysis in Section 3. As $\beta$ increases, the solution stabilizes more quickly and at a lower temperature. Although the model is simplified by using a two-dimensional domain and a uniform cooling term, it still shows clearly that stronger air conditioning can effectively reduce thermal risk in the kitchen. To make the model more realistic, we can step into the relationship between $\beta$ and air condition power, which will give a quantified standard for the AC, rather than representing the scale of AC by the abstract parameter $\beta$.

\section{Conclusion}
In this project, we developed a mathematical model to investigate the temperature distribution in a kitchen heated by a stove. The kitchen was represented as a two-dimensional square domain, and the heat transfer process was described by a PDE incorporating diffusion, a heat source from the stove, and a uniform cooling sink due to air conditioning. Numerical simulations were performed using the Backward Euler scheme to analyze how air-conditioning strength affects the thermal environment. Our results show that the temperature distributions evolve over time to reach a steady state, where the heat input from the stove is balanced by the cooling from the air conditioning.

\vspace{3mm}

\noindent By comparing different air-conditioning strengths, we demonstrated that stronger cooling reduces both the peak temperature experienced by the chef and the spatial extent of the high-temperature region. These results highlight how the balance between heat generation and dissipation influences the thermal safety of a confined workspace. Moreover, to make the chef work in a safe environment, the value of $\beta$ should be larger then a minimum value of $3.51 \ hour^{-1}$, in other words, the choice of air condition should be strong enough to ensure the  kitchen safety.

%%%%%%%%%%%%%%%%%%%% Refrence Start %%%%%%%%%%%%%%%%%%%%%%%%%%%
\bibliographystyle{plain} % you can choose from many different styles: https://www.overleaf.com/learn/latex/Bibtex_bibliography_styles
\begin{thebibliography}{4} % Entries are in the Refs.bib file
\bibitem[1]{one}
Thermtest Inc. 
Materials thermal properties database. \emph{Thermtest Thermal Resources}, 2023. Available at: https://thermtest.com/thermal-resources/materials-database.
\bibitem[2]{two}
Government of Canada.
Thermal stress in the work place.  \emph{Employment and Social Development Canada}, 2018. Available at: https://www.canada.ca/en/employment-social-development/services/health-safety/reports/thermal-stress-work-place.html.
% \bibitem[3]{three}
% W.K. Chow and Xiaomin Ni.
% Experimental evaluation on performance of open kitchen fire suppression systems. \emph{Fire Safety Science—Proceedings of the Eleventh International Symposium}, pages 1298–1311, 2014.
\bibitem[3]{three}
Dingming Zheng, Lei Su, Haoyu Ou, and Shijie Ruan.
Study on heat transfer characteristics and performance of the full premixed cauldron stove with porous media. \emph{Energies}, 15:9523, 2022.
\bibitem[4]{four}
Miranda Holmes-Cerfon. PDE-Lecture2-problems.pdf. MATH 461, University of British Columbia. 2026.
\end{thebibliography}
%%%%%%%%%%%%%%%%%%%% Reference end %%%%%%%%%%%%%%%%%%%%%%%%%%%%%

\newpage

%%%%%%%%%%%%%%%%%%%% Gen AI Declaration Start %%%%%%%%%%%%%%%%%%%%%%%%%
\section*{Gen AI declaration} 
Generative AI tools including ChatGPT and Grok were used as study aids for this project. They were used to help clarify thermodynamic and heat transfer concepts related to our model and locate relevant research articles on thermal properties and heat transfer.

\vspace{5mm}

\noindent All mathematical models, parameters, analysis, simulations, figures, and conclusions were developed and verified independently by the authors. All equations, assumptions, and results presented in this report reflect our own understanding. The final report was reviewed and edited by the authors to ensure accuracy, precision, and clarity.
%%%%%%%%%%%%%%%%%%%% Gen AI Declaration end %%%%%%%%%%%%%%%%%%%%%%%%%%%

\end{document}
